\section{Cubic pencils in genus five}\label{app:pencils}

Fix a smooth genus-five curve $D$. A genus-five cubic can move in a
pencil even after quotienting by dilation. At a cusp its normalization
has one fewer handle than the general member. Both genera occur in the
universal comparison.

\begin{theorem}\label{thm:pencil}
Let $N$ be a general stable bundle of rank two and determinant $K_D$.
In $\Tot N$ the reduced space of degree-three cycles is the union of
$6480$ lines and $32$ planes meeting only at $3[D]$. Dilation has
weight three on their parameters. A general determinant-one quadratic
shear along $D$ isolates $3[D]$ in a threefold $X$. For its unit $J_3$,
\[
 J_3-x_{5,3}=2160x_{16,1}+\tfrac{1280}{3}x_{18,1}
                            -\tfrac{32}{9}x_{19,1}.
\]
The difference of the connected fixed-target degree-three series in
the same order is $2160w^{15}+(1280/3)w^{17}-(32/9)w^{18}$.
\end{theorem}

\subsection{The pencils and their limits}
By Section~\ref{sec:evaluations} all nonzero cubics are $y^3=s$ on
saturated sublines of degrees one and two. For degree one
Corollary~\ref{cor:cubic-fiber-count} gives $6480$, with unique
reduced divisors by the dimension argument in Section~\ref{sec:covers}.
Each smooth cubic has genus 16. There are $2^5$ maximal degree-two
sublines, reduced for general $N$ \cite[Section~4]{Gillam}. One may
fix the determinant here since tensoring by a line reaches every
determinant of the same degree. At a degree-two subline $S$ the
projective extensions of $Q=K_DS^{-1}$ by $S$ have dimension seven,
so a proper closed locus of $S$ has image of dimension at most eleven
in the twelve-dimensional moduli of bundles. Thus $S^3$ is a general
degree-six line. Its two sections form a base-point-free pencil:
specialness requires $S^3=K_D(-B_2)$, a base point requires
$S^3=K_D(p-B_3)$, and a triple zero or two double zeros require
$S^3=\cO(3p+B_3)$ or $\cO(2p+2q+B_2)$, where $B_i$ is effective
of degree $i$. Each locus has dimension at most four. The pencil
therefore has twenty simple branch values by Riemann--Hurwitz. A
general cubic has six totally ramified points and genus 19. Each
exceptional member has one cusp $x^2=y^3$ and four other simple
zeros, so its normalization has genus 18. Saturation and the spectral
equation recover $(S,s)$ from the cycle; two general evaluations
bound all coefficients on each of the finitely many spaces. Hence
the bijection from their union is proper and gives its topology.

For the embedded cubic $F$ in $S\subset N$ the boundary in its
normal sequence is multiplication by $3y^2 e$, where
$e\in H^1(SQ^{-1})$ is the extension class. For a degree-two $S$
the nonzero connecting block is $H^0(Q)\to H^1(S)$, invertible by
acyclicity of $N$. Here $H^0(S^2)=0$ and $H^1(S^3)=0$. The finite
flat algebra $\cO\oplus S^{-1}\oplus S^{-2}$ also applies to the
embedded cusp. Thus throughout the punctured plane
\[
 H^0(N_{F/\Tot N})=H^0(S^3),\qquad H^1(N_{F/\Tot N})=
 H^1(QS^{-2})=H^0(S^3)^* .
\]
For a degree-one line $R$ the two additional connecting blocks are
dual to the pairing $(a,b)\mapsto\langle e,sab\rangle$ on
$H^0(K_DR^{-2})\times H^0(K_DR^{-3})$. Both spaces have dimension
two. Two general points off $\operatorname{div}s$ give independent
evaluations on both pencils, so a sum of their weighted evaluations
makes this form nondegenerate. Bad forms occupy a proper subset of
each extension fiber and the finite incidence excludes them for
general $N$. The remaining tangent and obstruction each have
dimension one.

Choose a determinant-one shear as in Section~\ref{sec:neighborhood}.
Its quadratic term on $S^2$, projected to $Q$, gives a fixed functional
in $H^0(S^3)^*$ over the entire punctured plane. On the normalization
of a cusp over $p$ the covering algebra has last eigenline
$S^{-2}(p)$; its obstruction is the restriction of the functional
to $H^0(S^3(-p))$. This is the line generated by the cusp section.
Principal parts span $H^1(QS^{-2})$, so finitely many open conditions
give a shear nonzero on all $640$ cusp lines, all 32 plane obstruction
spaces and all $6480$ isolated-line obstructions. At a smooth
embedded limit the relative Hilbert tangent space now has precisely
the coefficient directions. The actual coefficient space has that
dimension and hence fills its regular germ in the central target.

At a cuspidal limit the Hilbert tangent test alone would miss arcs
whose generic curves have smaller arithmetic genus. Suppose such a
noncentral arc exists. Curve selection in a projective Chow space
and properness of stable maps extend the normalizations over a DVR,
after finite base change. Write $t=cu^k+O(u^{k+1})$ for the target
parameter. The sole positive special component is birational to the
cusp; its normalization has the above eigenline decomposition.
Since $H^0(K_DS^{-3}(p))=0$ by Riemann--Roch and
$h^0(S^3(-p))=1$, and $H^0(K_DS^{-2})=0$ by $2$-rigidity and
Riemann--Roch, all sections of the pulled-back $Q$ on that component
come from $D$. Matching gives the same assertion on contracted pieces.
Over $A=\CC[[u]]/(u^k)$ all sections of $g^*Q$ for the map
$g:\mathcal C_A\to D_A$ still come from $D_A$: subtract the lift of
each successive special-fiber section and use source flatness.
The target is linear over $A$. The normal height in $g^*N$ has
zero trace since $H^0(D,N)=0$. This trace is defined even with
contracted components: $Rg_*\cO_{\mathcal C_A}$ is perfect by the
regular graph into $\mathcal C_A\times_A D_A$ and proper flat
perfect projection. Its rank on $D_A$ is three. Hence a pulled-back
quotient height has trace three times its base value, so the quotient
height vanishes modulo $u^k$. Modulo $u^{k+1}$ its first possible
coefficient would have to trivialize the projected quadratic
cocycle times $c$ on the special source. Its restriction to the
normalization is nonzero in $H^1(QS^{-2}(p))$, a contradiction.
This excludes noncentral arcs. Rescaling any sequence of displaced
cubic cycles to a fixed height, as in Proposition~\ref{prop:isolation},
would give one of the smooth or cuspidal limits just excluded.
Thus the sheared triple is isolated.

\subsection{The normalization of one cusp}
Near a nonzero cuspidal member the connected genus-18 map is its
smooth normalization, with trivial stabilizer. Any node, positive-genus
contracted component or graph loop would raise the arithmetic genus.
A contracted rational tree with one attachment would be unstable.
On an Artin base the map to $D$ stays finite flat. The same trace
and projected obstruction, now applied to consecutive powers of the
maximal ideal, force $t=0$ and the height into $S$. Its characteristic
equation is $y^3=s_A$ by $H^0(S)=H^0(S^2)=0$.
In local coordinates near the double zero the equation is
$X^2=Y^3+b$, with $b$ the critical value and $a$ the independent
parameter along the cusp ray. On the map the parametric Morse lemma
gives $Y=z^2+r$, with $r$ in the maximal ideal and the coordinate
sign fixed on the central fiber. The equality $X(z)^2=X(-z)^2$ and
regularity of $X(z)-X(-z)$ (reduce first to $2z^3$) give
$X=z h(z^2)$. Then
\[
 b=-r^3,\qquad h(v)^2=v^2+3rv+3r^2 .
\]
Weierstrass preparation forces $h=v+d$ with $2d=3r$ and $d^2=3r^2$.
Hence $r^2=b=0$. Conversely $(Y,X)=(z^2+r,z^3+3rz/2)$ for any
$r^2=0$ gives the required local normalization family. Its finite
free rank-three algebra agrees with the spectral algebra off the
marked disk by reduction and Nakayama. The two height functions
determine the gluing on the punctured disk uniquely. Thus the completed
map ring in a fiber frame is $\CC[[a,r]]/(r^2)$.
With two further smooth base directions this is an lci of the absolute
virtual dimension three. Its two-term obstruction theory has the same
rank as the lci cotangent complex. At any residue field the degree-zero
map is an isomorphism and the degree-minus-one map is surjective;
equality of ranks and derived Nakayama identify the two complexes.
The local virtual multiplicity is two.

\subsection{Moving the two evaluations}
Work over $B=\PP^2$ with tautological dual $M=\cO(1)$. Two general
fiber products give an injective map on each line, an isomorphism on
every plane and a Thom product representing $9x_{5,3}$. Translate by
two general small cubic forms in $M^3$. Move the shear parameter
graph to the metric dual $\tau$ of a linear section. Its zero line
$T$ has oriented normal index $-1$. The translations have nine common
zeros off $T$, accounting for $9J_3$. Their ratio on $T$ has degree
three, unramified over the finitely many line and cusp directions.
Each isolated line gives three negative genus-16 intersections.
Over each plane there remains a compact family above the whole of
$T$, with three incidences per cusp direction. At each of these the
two translations cut the amplitude and the ratio transversely, so the
genus-18 virtual degree is $-2$. Compactness of this homotopy follows
from the uniform lower bound on the pair of linear evaluations on
the unit link of the finite union of planes and lines, together with
isolation on the small noncentral fibers and the disk mass estimate.

To extend the cusp calculation to all source genera we deform each
whole pencil, since a genus-19 image can approach a cusp. Add a second
quadratic parameter spanning the two-dimensional embedded obstruction
with the first at every plane. On the graph take a second small
metric-dual linear section, cut off away from $T$ near the nine
marked points. Its zero on $T$ is a single general point. At this
point the two anti-complex normal factors have positive product,
yielding one positive smooth genus-19 intersection per plane.
For control of the other images take a relative projective completion.
The uniform mass bound bounds their ample degrees on the disk.
Nearby cycles at a nonzero plane member are integral with multiplicity
one, by proper Chow addition. In bounded degree their arithmetic
genera are bounded by the generic plane-projection genus. A Hilbert
limit with fundamental cycle $C$ surjects onto $\cO_C$ with a
finite-length kernel, giving $p_a(\mathrm{generic})\le p_a(C)$
with equality only for the reduced Hilbert limit. Stable-map limits
give the opposite lower bound on geometric genus by the normalization
of $C$. Thus near smooth members the arithmetic genus equals 19 and
the reduced Hilbert germ forces both parameters to vanish. Near a cusp
away from the common parameter zero the same argument excludes
arithmetic genus 19, while geometric genus is at least 18.
Every remaining small-parameter image there is therefore smooth
embedded of genus 18. Compactness gives uniform neighborhoods.

At these images the shears at fresh unramified sheets of
Lemma~\ref{lem:shear-jets} make the Hilbert chart regular; one may use
powers two through four of the separating covector, whose values are
nonzero and distinct, to preserve the first jet at the core. The new
three columns interpolate on the three sheets by the Vandermonde
determinant. A small smooth change of graph and translations gives
transverse isolated curves. Excision in source genus 18 preserves the
degree $-2$ in each cusp box. Proposition~\ref{prop:moving-embedded-curve}
now assigns the entire series $w^{17}$ to each signed genus-18 image,
including all contracted components. The same signed multiplicity-one
units arise by proper connected pullback. The complete unnormalized
cut is consequently
\[
 9J_3-3\cdot6480x_{16,1}-2\cdot60\cdot32x_{18,1}+32x_{19,1}.
\]
Dividing by nine gives the two identities in Theorem~\ref{thm:pencil}.
For the complex class, the two lower degrees agree by the empty-tuple
argument; hence the difference after subtraction is primitive and
killed by proper division. Its smooth image vanishes by the signed
calculation, so it vanishes by multiplicity-one detection.
Independently the handle recurrence gives
$9D_4\equiv3w^{17}+5w^{18}\pmod9$, consistent with the displayed signs.
